\usepackage{amsmath}
\usepackage{amssymb}
\usepackage[spanish]{babel}

\title{TP LaTeX 3Compu, un infiltrado}
\author{
    \textbf{Integrantes:} \\
        Hector Gomez\\
        Leonardo Bravo\\
        Avril Tapia\\ 
        Lautaro Garcia
    }
\date{Septiembre 13 2026}

% comandos personalizados !

% Math
\newcommand{\N}{\mathbb{N}}
\newcommand{\Z}{\mathbb{Z}}
\newcommand{\R}{\mathbb{R}}

% math secuencias
\newcommand{\seq}[1]{\langle#1\rangle}
\newcommand{\seqq}{\subseteq}
\newcommand{\set}[3]{\text{setAt} (#1,#2,#3)}



% math diccionarios
\newcommand{\dict}[1]{\text{dict}\{#1\}}


% espacios
\newcommand{\QN}[1]{\hspace*{#1em}}

% implicaciones
\newcommand{\imp}{\rightarrow}
\newcommand{\impl}{\rightarrow_{\text{L}}}

% types
% #1 = nombre del tipo
% #2 = descripcion del tipo
\newcommand{\type}[2]{&\text{type } #1:~#2\\}

% observadores
% #1 = nombre del observador
% #2 = tipado del observador
\newcommand{\obs}[2]{&\QN2\text{obs } #1:~#2\\}

% func aux
% #1 = nombre de la funcion
% #2 = funcion
\newcommand{\aux}[2]{&\QN2\textbf{aux } \text{#1} = #2}

% requiere y asegura

% requiere
% #1 = condiciones de requiere
\newcommand{\req}[1]{%
  & \QN{4} \textbf{requiere:}~\{ \\%
  & \QN{6} #1 \\%
  & \QN{4} \}%
}

% asegura
% #1 = condiciones de asegura
\newcommand{\aseg}[1]{%
  & \QN{4} \textbf{asegura:}~\{ \\%
  & \QN{6} #1 \\%
  & \QN{4} \}%
}

% proc 
% #1 = nombre del procedimiento
% #2 = parametros de entrada
% #3 = parametros de salida
% #4 = requiere
% #5 = asegura
\newcommand{\proc}[5]{%
  & \QN{2} \textbf{proc } \text{#1} (#2): \text{#3} \{ \\%
  \req{#4} \\%
  \aseg{#5} \\%
  & \QN{2} \}%
}

% tad
% #1 = nombre del TAD
% #2 = contenido del TAD (procedimientos, auxiliares, etc)
\newcommand{\tad}[2]{%
  \begin{aligned}
    & \textbf{TAD: } \text{#1} \{ \\
    #2 \\
    & \}
  \end{aligned}%
}

% IDEA PRINCIPAL:
% Imp un solo \tad {nombre} {contenido}
% contenido iremos poniendo \proc{}{}{}{}{}\\
%                           \aux{}{}\\
% y si tenemos que hacer pred, hacemos otro command para eso y lo agregamos
